\chapter{Kết luận}
\label{Chapter6}

\section{Ý nghĩa kết quả đạt được}
\par Với đề tài \textit{``Xây dựng hệ thống theo dõi bệnh nhân từ xa cho người bệnh tăng huyết áp và đái tháo đường''}, mục tiêu ban đầu của nhóm là xây dựng một hệ thống hỗ trợ bệnh nhân ghi nhận các chỉ số sức khỏe hằng ngày, đồng thời giúp bác sĩ và nhân viên y tế có thể theo dõi tình trạng bệnh nhân từ xa thông qua các nền tảng số. Đây là một bài toán có ý nghĩa thực tiễn vì các bệnh mãn tính như tăng huyết áp và đái tháo đường thường yêu cầu quá trình theo dõi lâu dài, liên tục và có tính lặp lại. Nếu chỉ dựa vào các lần tái khám định kỳ tại cơ sở y tế, bác sĩ rất khó nắm bắt đầy đủ diễn biến sức khỏe của bệnh nhân trong đời sống hằng ngày. Vì vậy, hệ thống mà nhóm xây dựng hướng đến việc tạo ra một kênh thu thập dữ liệu sức khỏe có tổ chức, giúp quá trình theo dõi trở nên thuận tiện và chủ động hơn.

\par Nhìn chung, nhóm đã hoàn thành được phần lớn các mục tiêu cốt lõi đã đề ra trong phạm vi đồ án. Ứng dụng di động dành cho bệnh nhân đã hỗ trợ các chức năng quan trọng như đăng nhập, đăng ký, nhập chỉ số sức khỏe, xem lịch sử đo, theo dõi cảnh báo, nhận thông báo, trao đổi tin nhắn với bác sĩ và quản lý hồ sơ cá nhân. Bên cạnh đó, ứng dụng di động cũng hỗ trợ luồng sử dụng dành cho y tá, cho phép y tá xem danh sách bệnh nhân được phân công và nhập chỉ số sức khỏe thay cho bệnh nhân trong một số trường hợp cần thiết. Việc bổ sung vai trò y tá giúp hệ thống phản ánh tốt hơn bối cảnh vận hành thực tế tại cơ sở y tế, nơi bệnh nhân không phải lúc nào cũng tự nhập dữ liệu.

\par Đối với phía bác sĩ, nhóm đã triển khai được ứng dụng web giúp bác sĩ theo dõi danh sách bệnh nhân được phân công, xem thông tin chi tiết bệnh nhân, quan sát lịch sử đo và biểu đồ xu hướng, cấu hình ngưỡng an toàn, tiếp nhận cảnh báo khi chỉ số vượt ngưỡng, tạo nhắc nhở và trao đổi với bệnh nhân thông qua chức năng chat. Những chức năng này giúp bác sĩ có thêm công cụ hỗ trợ trong quá trình theo dõi bệnh nhân từ xa, đặc biệt là đối với các chỉ số cần được ghi nhận thường xuyên như huyết áp, đường huyết, nhịp tim, SpO2, thân nhiệt và nhịp thở. Thay vì chỉ xem dữ liệu tại thời điểm tái khám, bác sĩ có thể quan sát dữ liệu theo chuỗi thời gian, từ đó nhận biết xu hướng thay đổi rõ ràng hơn.

\par Ngoài các thành phần dành cho bệnh nhân, y tá và bác sĩ, nhóm cũng đã xây dựng trang web dành cho quản trị viên nhằm hỗ trợ vận hành hệ thống. Trang quản trị cho phép quản lý tài khoản bác sĩ, bệnh nhân, y tá, khoa phòng, phân công nhân viên y tế phụ trách bệnh nhân và theo dõi lịch sử hoạt động. Việc có riêng một trang quản trị giúp hệ thống phân tách rõ ràng giữa nghiệp vụ chuyên môn của bác sĩ và nghiệp vụ vận hành của quản trị viên. Điều này làm cho mô hình hệ thống đầy đủ hơn, đồng thời tạo nền tảng để mở rộng trong tương lai nếu triển khai cho nhiều khoa phòng hoặc nhiều nhóm bệnh nhân khác nhau.

\par Về mặt kỹ thuật, nhóm đã xây dựng được một hệ thống nhiều thành phần, bao gồm ứng dụng di động, ứng dụng web, trang quản trị, backend, cơ sở dữ liệu, cơ chế xử lý nền, thông báo đẩy và hạ tầng triển khai. Backend được xây dựng bằng Go và Gin, sử dụng MongoDB để lưu trữ dữ liệu, Redis để hỗ trợ một số luồng thời gian thực, Temporal để xử lý các workflow nền, Firebase Cloud Messaging để gửi thông báo đẩy và WebSocket để hỗ trợ chat hoặc cập nhật sự kiện theo thời gian gần thực. Hệ thống cũng được đóng gói bằng Docker và triển khai trên AWS, trong đó backend chạy trên Amazon EC2 còn các ứng dụng web được lưu trữ dưới dạng static files trên Amazon S3 và phục vụ người dùng qua Cloudflare để thiết lập HTTPS.

\par Kết quả đạt được của đề tài không nhằm thay thế hoạt động khám chữa bệnh trực tiếp hoặc đưa ra chẩn đoán y khoa tự động, mà đóng vai trò như một công cụ hỗ trợ theo dõi và quản lý dữ liệu sức khỏe. Hệ thống giúp bệnh nhân chủ động hơn trong việc ghi nhận chỉ số, giúp bác sĩ có thêm dữ liệu tham khảo và giúp nhân viên y tế phát hiện sớm một số dấu hiệu bất thường dựa trên cơ chế cảnh báo theo ngưỡng. Đây là nền tảng quan trọng để tiếp tục phát triển các hệ thống theo dõi sức khỏe thông minh hơn trong tương lai.

\section{Kiến thức, kinh nghiệm rút ra từ đề tài}
\par Trong quá trình thực hiện đề tài, nhóm đã đối mặt với nhiều khó khăn liên quan đến phân tích nghiệp vụ, thiết kế hệ thống, xây dựng giao diện, phát triển backend, tích hợp giữa các thành phần, xử lý thông báo, triển khai cloud và kiểm thử theo nhiều vai trò người dùng. Tuy nhiên, thông qua quá trình tự học, trao đổi nhóm, thử nghiệm và sửa lỗi liên tục, các thành viên đã tích lũy được nhiều kiến thức và kinh nghiệm có giá trị. Những kinh nghiệm này không chỉ phục vụ cho việc hoàn thành đồ án, mà còn có thể áp dụng trong các dự án phần mềm thực tế sau này.

\par Một trong những bài học quan trọng nhất mà nhóm rút ra là việc xây dựng một hệ thống phần mềm hoàn chỉnh không chỉ dừng lại ở việc cài đặt các màn hình hoặc API riêng lẻ. Đối với hệ thống RPM, mỗi chức năng thường liên quan đến nhiều thành phần khác nhau. Ví dụ, khi bệnh nhân nhập một chỉ số sức khỏe, ứng dụng di động phải kiểm tra dữ liệu và gửi request đúng định dạng, backend phải xác thực người dùng và lưu dữ liệu, hệ thống xử lý nền phải đánh giá dữ liệu với ngưỡng an toàn, cảnh báo và thông báo phải được tạo nếu có vi phạm, còn bác sĩ cần thấy dữ liệu mới trên giao diện web. Vì vậy, nhóm cần thiết kế hệ thống theo luồng end-to-end và kiểm thử theo cách người dùng thật sẽ sử dụng sản phẩm.

\subsection{Về mặt kỹ thuật}

\subsubsection{Phát triển hệ thống}
\par Trong quá trình phát triển các thành phần của hệ thống, nhóm đã học được nhiều kiến thức kỹ thuật quan trọng, từ phát triển giao diện người dùng, xây dựng backend, thiết kế cơ sở dữ liệu, xử lý workflow nền đến triển khai hệ thống trên cloud. Những kiến thức này giúp nhóm hiểu rõ hơn cách một hệ thống phần mềm nhiều thành phần được tổ chức và vận hành.

\begin{itemize}
    \item Frontend web: Nhóm đã sử dụng React, TypeScript và Vite để xây dựng ứng dụng web dành cho bác sĩ và trang quản trị. Thông qua quá trình phát triển, nhóm học được cách tổ chức mã nguồn theo component, chia nhỏ giao diện thành các thành phần có thể tái sử dụng, quản lý route bằng React Router, xử lý trạng thái đăng nhập, tích hợp API bằng Axios và hiển thị dữ liệu dưới dạng bảng hoặc biểu đồ. Đặc biệt, các màn hình như dashboard, danh sách bệnh nhân, chi tiết bệnh nhân, cấu hình ngưỡng, cảnh báo, nhắc nhở và quản lý người dùng giúp nhóm có thêm kinh nghiệm trong việc xây dựng các giao diện quản lý dữ liệu tương đối phức tạp.

    \item Ứng dụng di động: Nhóm đã sử dụng Expo và React Native để phát triển ứng dụng di động dành cho bệnh nhân và y tá. Trong quá trình thực hiện, nhóm học được cách xây dựng giao diện mobile đa màn hình, điều hướng theo vai trò người dùng, lưu trữ thông tin phiên đăng nhập, gọi API từ thiết bị di động, xử lý form nhập liệu, hiển thị lịch sử đo, cảnh báo, thông báo và tin nhắn. Việc phát triển mobile cũng giúp nhóm hiểu rõ hơn các vấn đề đặc thù của thiết bị di động như kích thước màn hình khác nhau, bàn phím che giao diện, quyền nhận thông báo, token thiết bị và kiểm thử trên thiết bị thật hoặc giả lập.

    \item Backend: Nhóm đã sử dụng Go và Gin để xây dựng hệ thống backend theo hướng tách biệt giữa router, handler, service, repository, middleware và domain model. Quá trình này giúp nhóm hiểu rõ hơn về cách xây dựng RESTful API, cách tổ chức mã nguồn backend có khả năng bảo trì, cách xử lý request/response, cách xác thực bằng JWT, cách kiểm tra vai trò người dùng và cách kết nối backend với các dịch vụ khác. Việc sử dụng Go cũng giúp nhóm làm quen với cách quản lý package, cách xử lý lỗi, cách tổ chức dependency và cách build ứng dụng thành binary hoặc Docker image.

    \item Cơ sở dữ liệu: Nhóm đã sử dụng MongoDB làm cơ sở dữ liệu chính cho hệ thống. Thông qua việc thiết kế các collection như users, measurements, thresholds, alerts, reminders, assignments, departments, conversations, messages, notifications và activity logs, nhóm học được cách tổ chức dữ liệu dạng document, cách liên kết dữ liệu thông qua ObjectId, cách thiết kế dữ liệu linh hoạt nhưng vẫn đảm bảo phục vụ được nghiệp vụ. Đối với hệ thống có nhiều vai trò và nhiều luồng dữ liệu như RPM, việc thiết kế cơ sở dữ liệu hợp lý có ảnh hưởng rất lớn đến khả năng mở rộng và khả năng truy vấn dữ liệu.

    \item Xử lý nền và workflow: Việc tích hợp Temporal giúp nhóm hiểu rõ hơn về cách xử lý các tác vụ nền không nên đặt trực tiếp trong HTTP request. Các luồng như đánh giá dữ liệu đo để sinh cảnh báo hoặc gửi nhắc nhở theo lịch đều phù hợp để triển khai bằng workflow. Thông qua quá trình này, nhóm học được cách tách logic xử lý dài hoặc lặp lại ra khỏi API server, cách định nghĩa workflow, activity và cách kiểm tra kết quả xử lý nền. Đây là kinh nghiệm quan trọng vì nhiều hệ thống thực tế cần xử lý các tác vụ bất đồng bộ, retry hoặc thực thi theo lịch.

    \item Thời gian thực và thông báo: Nhóm đã tích hợp WebSocket, Redis Pub/Sub và Firebase Cloud Messaging để phục vụ các nhu cầu cập nhật dữ liệu theo thời gian gần thực. Thông qua các chức năng như chat giữa bệnh nhân với bác sĩ, thông báo cảnh báo mới và push notification, nhóm học được rằng việc truyền tải sự kiện thời gian thực phức tạp hơn so với các API thông thường. Hệ thống cần quản lý kết nối, token thiết bị, trạng thái gửi thông báo, trường hợp người dùng offline và cách đồng bộ dữ liệu giữa nhiều client khác nhau.

    \item Triển khai hệ thống: Nhóm đã sử dụng Docker, Docker Compose, Nginx, GitHub Actions, AWS và Cloudflare để triển khai hệ thống. Backend chạy trên Amazon EC2, các ứng dụng web được lưu trữ trên Amazon S3 và phục vụ người dùng qua Cloudflare để thiết lập HTTPS, còn quy trình CI/CD hỗ trợ tự động build và cập nhật hệ thống khi mã nguồn thay đổi. Quá trình triển khai giúp nhóm học được cách cấu hình biến môi trường, reverse proxy, HTTPS, CORS, WebSocket, security group, Docker network và kiểm tra log container.
\end{itemize}

\subsubsection{Sử dụng thuần thục các công cụ hỗ trợ}
\par Bên cạnh kiến thức lập trình, nhóm cũng tích lũy được nhiều kinh nghiệm khi sử dụng các công cụ hỗ trợ phát triển phần mềm. Việc sử dụng các công cụ này một cách có hệ thống giúp quá trình làm việc nhóm rõ ràng, minh bạch và hiệu quả hơn.

\begin{itemize}
    \item GitHub Organization: Nhóm sử dụng GitHub Organization để quản lý mã nguồn tập trung cho các thành phần của hệ thống. Thông qua quá trình làm việc, các thành viên rèn luyện thêm kỹ năng tạo branch, đặt tên branch, commit theo từng nhóm thay đổi, tạo pull request, review mã nguồn, giải quyết conflict và bảo vệ nhánh chính. Đây là quy trình rất cần thiết trong các dự án có nhiều người cùng phát triển.

    \item Jira Software: Nhóm sử dụng Jira để quản lý công việc, chia sprint, tạo backlog, phân công ticket và theo dõi tiến độ. Việc sử dụng Jira giúp nhóm học được cách mô tả công việc rõ ràng, ước lượng khối lượng, xác định trạng thái công việc và theo dõi những phần còn tồn đọng. Nhờ đó, nhóm hạn chế được tình trạng quên việc, giao việc bằng miệng hoặc không biết chức năng nào đang do ai phụ trách.

    \item Google Drive: Nhóm sử dụng Google Drive để lưu trữ tài liệu chung như đề cương, báo cáo, biên bản họp, sơ đồ, hình ảnh, tài liệu API và các file hỗ trợ khác. Việc lưu trữ tập trung giúp các thành viên dễ dàng truy cập tài liệu mới nhất, hạn chế việc mỗi người giữ một phiên bản riêng và giảm rủi ro mất dữ liệu.

    \item Google Meet: Nhóm sử dụng Google Meet để họp định kỳ vào thứ 7 hằng tuần và tổ chức các buổi họp bổ sung khi có vấn đề quan trọng. Thông qua các buổi họp, nhóm có thể báo cáo tiến độ, demo chức năng, chia sẻ màn hình để debug, thảo luận giải pháp và thống nhất công việc cho tuần tiếp theo. Điều này giúp cải thiện sự phối hợp giữa các thành viên, đặc biệt khi các phần backend, frontend và mobile có nhiều điểm phụ thuộc lẫn nhau.

    \item diagrams.net: Nhóm sử dụng diagrams.net để vẽ các sơ đồ kiến trúc hệ thống, sơ đồ triển khai AWS, sơ đồ cơ sở dữ liệu, sơ đồ workflow cảnh báo, workflow nhắc nhở và các sơ đồ luồng xử lý khác. Việc trực quan hóa hệ thống giúp các thành viên dễ thống nhất cách hiểu, đồng thời giúp báo cáo có tính rõ ràng và chuyên nghiệp hơn.

    \item Figma: Nhóm sử dụng Figma để thiết kế giao diện và prototype cho ứng dụng mobile, web bác sĩ và trang quản trị. Việc thiết kế giao diện trước khi cài đặt giúp frontend và mobile có định hướng rõ ràng, giảm tình trạng vừa làm vừa đoán, đồng thời giúp cả nhóm góp ý sớm trước khi bước vào giai đoạn lập trình.
\end{itemize}

\subsection{Về kỹ năng làm việc nhóm và lên kế hoạch}
\par Kỹ năng làm việc nhóm, phân chia công việc và lên kế hoạch có ảnh hưởng rất lớn đến kết quả của đề tài. Do hệ thống có nhiều thành phần được phát triển song song, nếu nhóm không có quy trình phối hợp rõ ràng thì rất dễ xảy ra tình trạng backend hoàn thành nhưng frontend chưa có giao diện, frontend cần dữ liệu nhưng API chưa sẵn sàng, hoặc mobile tích hợp sai định dạng response. Qua quá trình thực hiện, nhóm đã rút ra nhiều kinh nghiệm quan trọng trong cách tổ chức công việc.

\begin{itemize}
    \item Cần phân chia nhiệm vụ dựa trên thế mạnh của từng thành viên. Trong đề tài này, Nhân đảm nhiệm chính phần backend và hosting, Hiển đảm nhiệm backend và mobile, còn Khoa và Kiệt tập trung vào frontend. Cách phân chia theo mảng kỹ thuật giúp mỗi thành viên có thể tập trung phát triển chuyên sâu, đồng thời giảm sự chồng chéo không cần thiết trong quá trình làm việc.

    \item Cần mô tả yêu cầu thật rõ trước khi bắt đầu cài đặt. Một ticket không nên chỉ ghi chung chung như ``làm màn hình cảnh báo'' hoặc ``viết API đo chỉ số'', mà cần mô tả người dùng nào sử dụng, dữ liệu đầu vào là gì, dữ liệu đầu ra là gì, endpoint nào được gọi, giao diện cần hiển thị những trạng thái nào và tiêu chí hoàn thành là gì. Việc mô tả rõ yêu cầu giúp giảm số lần phải sửa lại do hiểu sai.

    \item Cần thống nhất API contract trước khi frontend và mobile tích hợp. Đối với các chức năng có sự tham gia của nhiều thành phần, backend nên mô tả trước method, endpoint, request body, response mẫu và các mã lỗi có thể xảy ra. Frontend và mobile có thể dựa vào đó để xây dựng giao diện hoặc dùng dữ liệu giả trước, giúp quá trình phát triển song song diễn ra hiệu quả hơn.

    \item Cần kiểm thử theo luồng hoàn chỉnh thay vì chỉ kiểm thử riêng từng phần. Một chức năng chỉ nên được xem là hoàn thành khi nó hoạt động đúng từ phía người dùng đến phía server và quay lại giao diện. Ví dụ, chức năng cảnh báo cần được kiểm tra từ lúc bệnh nhân nhập chỉ số, backend lưu dữ liệu, workflow tạo cảnh báo, thông báo được gửi và bác sĩ nhìn thấy cảnh báo trên web.

    \item Cần duy trì họp nhóm định kỳ để nắm tiến độ và giải quyết vướng mắc. Lịch họp vào thứ 7 hằng tuần giúp nhóm có thời điểm cố định để tổng kết sprint, demo chức năng, báo cáo khó khăn và lên kế hoạch cho tuần tiếp theo. Việc họp đều đặn giúp nhóm phát hiện sớm các phần bị chậm tiến độ hoặc có nguy cơ ảnh hưởng đến các thành viên khác.

    \item Khi xảy ra lỗi nghiêm trọng, điều quan trọng nhất là xác định nguyên nhân và khắc phục thay vì tìm người chịu trách nhiệm. Đặc biệt với các lỗi triển khai hoặc lỗi tích hợp nhiều thành phần, nguyên nhân có thể nằm ở frontend, backend, biến môi trường, reverse proxy, database hoặc mạng. Do đó, nhóm cần phối hợp kiểm tra log và dữ liệu theo từng bước để tìm đúng nguyên nhân.
\end{itemize}

\subsection{Một số kiến thức và kinh nghiệm khác}

\subsubsection{Kỹ năng thiết kế giao diện người dùng}
\par Trong quá trình thiết kế giao diện, nhóm nhận thấy rằng giao diện của một hệ thống y tế cần ưu tiên tính rõ ràng, dễ sử dụng và hạn chế gây nhầm lẫn. Đối tượng sử dụng của hệ thống không chỉ là người trẻ quen dùng công nghệ, mà còn có thể bao gồm bệnh nhân lớn tuổi, người mắc bệnh mãn tính hoặc người cần thao tác nhanh để nhập dữ liệu sức khỏe hằng ngày. Vì vậy, các màn hình mobile cần có bố cục đơn giản, nút bấm rõ ràng, nhãn chỉ số dễ hiểu, đơn vị đo đầy đủ và thông báo lỗi cụ thể khi người dùng nhập sai dữ liệu.

\par Đối với ứng dụng web dành cho bác sĩ, thách thức lớn hơn nằm ở việc hiển thị nhiều dữ liệu y tế trên cùng một giao diện mà vẫn đảm bảo dễ theo dõi. Bác sĩ cần xem thông tin bệnh nhân, lịch sử đo, biểu đồ xu hướng, cảnh báo, ngưỡng an toàn, nhắc nhở và tin nhắn. Nếu giao diện không được tổ chức tốt, người dùng sẽ khó tìm thông tin quan trọng hoặc mất nhiều thời gian thao tác. Do đó, nhóm học được cách chia giao diện thành các khu vực chức năng rõ ràng, sử dụng bảng, biểu đồ, bộ lọc, thẻ thống kê và modal để trình bày dữ liệu một cách hợp lý.

\par Đối với trang quản trị, nhóm rút ra kinh nghiệm rằng các màn hình quản lý cần nhất quán về cách trình bày bảng, form, nút thao tác, modal xác nhận và thông báo kết quả. Các chức năng như quản lý bác sĩ, bệnh nhân, y tá, khoa phòng và phân công có nhiều điểm tương đồng, vì vậy việc tái sử dụng component và thống nhất cách hiển thị giúp giảm thời gian phát triển, đồng thời giúp quản trị viên dễ làm quen hơn khi chuyển giữa các màn hình.

\subsubsection{Triển khai và vận hành hệ thống}
\par Thông qua đề tài, nhóm đã có thêm nhiều kinh nghiệm trong việc triển khai và vận hành một hệ thống phần mềm trên môi trường cloud. Khi chạy ở local, nhiều thành phần có thể hoạt động bình thường vì địa chỉ server, biến môi trường và cấu hình mạng đơn giản. Tuy nhiên, khi đưa lên AWS, nhóm phải xử lý nhiều vấn đề thực tế hơn như cấu hình EC2, security group, Docker Compose, Nginx reverse proxy, domain API, HTTPS, CORS, WebSocket, biến môi trường production và kết nối đến các dịch vụ bên ngoài.

\par Nhóm cũng nhận thấy việc quản lý thông tin nhạy cảm là vấn đề rất quan trọng. Các thông tin như chuỗi kết nối cơ sở dữ liệu, JWT secret, cấu hình Firebase, Cloudinary, OAuth, token hoặc cookie không nên được đưa trực tiếp vào mã nguồn. Trong môi trường triển khai, các thông tin này cần được quản lý thông qua biến môi trường, GitHub Secrets hoặc file cấu hình riêng trên server. Đây là một bài học quan trọng vì hệ thống xử lý thông tin cá nhân và dữ liệu sức khỏe, vốn là các loại dữ liệu cần được bảo vệ cẩn thận.

\par Ngoài ra, quá trình deploy cũng giúp nhóm hiểu rõ hơn vai trò của logging và monitoring. Khi hệ thống có lỗi ở môi trường production, việc chỉ nhìn giao diện người dùng là không đủ để xác định nguyên nhân. Nhóm cần kiểm tra log của container, log của API server, log Nginx, trạng thái Docker Compose, kết nối database, trạng thái Redis, Temporal và response trả về từ các dịch vụ bên ngoài. Kinh nghiệm này giúp nhóm có cái nhìn thực tế hơn về quá trình vận hành một hệ thống nhiều thành phần.

\subsubsection{Làm việc với dữ liệu sức khỏe và nghiệp vụ y tế}
\par Một bài học quan trọng khác là dữ liệu sức khỏe cần được xử lý thận trọng hơn so với nhiều loại dữ liệu thông thường. Các chỉ số như huyết áp, đường huyết, SpO2, nhịp tim hoặc thân nhiệt không chỉ là các con số hiển thị trên giao diện, mà có thể ảnh hưởng đến cách bệnh nhân và nhân viên y tế nhận định tình trạng sức khỏe. Do đó, hệ thống cần trình bày dữ liệu rõ ràng, tránh gây hiểu nhầm và luôn nhấn mạnh rằng cảnh báo của hệ thống chỉ mang tính hỗ trợ theo dõi, không phải là kết luận chẩn đoán hoặc chỉ định điều trị.

\par Nhóm cũng nhận thấy việc thiết kế ngưỡng cảnh báo cần có sự cân bằng. Nếu ngưỡng quá nhạy, hệ thống có thể tạo ra quá nhiều cảnh báo, khiến bác sĩ và y tá khó ưu tiên xử lý. Nếu ngưỡng quá rộng, hệ thống có thể bỏ sót các dấu hiệu bất thường đáng chú ý. Trong phạm vi đồ án, nhóm sử dụng cơ chế cảnh báo dựa trên ngưỡng do bác sĩ cấu hình hoặc ngưỡng cố định, vì cách tiếp cận này đơn giản, dễ giải thích và phù hợp với mô hình thử nghiệm. Tuy nhiên, nếu triển khai thực tế, việc xác định ngưỡng cần có sự tham gia của chuyên gia y tế và cần được kiểm chứng trên dữ liệu thật.

\section{Những vấn đề còn tồn đọng}
\par Mặc dù hệ thống đã hoàn thành các chức năng cốt lõi trong phạm vi đồ án, vẫn còn một số vấn đề cần được tiếp tục cải thiện nếu nhóm phát triển sản phẩm ở các giai đoạn tiếp theo. Những vấn đề này chủ yếu liên quan đến độ chính xác dữ liệu, khả năng mở rộng, bảo mật, trải nghiệm người dùng, kiểm thử thực tế và mức độ phù hợp với quy trình y tế thật.

\par Vấn đề đầu tiên là dữ liệu sức khỏe hiện vẫn được nhập thủ công bởi bệnh nhân hoặc y tá, chưa được đồng bộ trực tiếp từ các thiết bị đo như máy đo huyết áp, máy đo đường huyết, đồng hồ thông minh hoặc thiết bị theo dõi sức khỏe tại nhà. Cách nhập thủ công phù hợp với phạm vi Proof-of-Concept vì dễ triển khai, không phụ thuộc thiết bị và giảm chi phí ban đầu. Tuy nhiên, nó vẫn tồn tại rủi ro người dùng nhập sai, nhập thiếu, nhập không đúng thời điểm đo hoặc quên nhập dữ liệu. Điều này có thể ảnh hưởng đến độ tin cậy của lịch sử đo và kết quả cảnh báo.

\par Vấn đề thứ hai là cơ chế cảnh báo của hệ thống hiện dựa trên ngưỡng, chưa sử dụng các mô hình phân tích nâng cao hoặc học máy. Cách tiếp cận theo ngưỡng có ưu điểm là dễ hiểu và dễ kiểm soát, nhưng chưa thể phản ánh đầy đủ bối cảnh lâm sàng của từng bệnh nhân. Trong thực tế, việc đánh giá một chỉ số bất thường có thể phụ thuộc vào tuổi, bệnh nền, thuốc đang sử dụng, thời điểm đo, triệu chứng đi kèm, mức độ vận động trước khi đo và nhiều yếu tố khác. Do đó, cảnh báo hiện tại chỉ nên được xem như tín hiệu hỗ trợ để bác sĩ hoặc nhân viên y tế xem xét thêm, không phải là kết luận cuối cùng về tình trạng sức khỏe.

\par Vấn đề thứ ba là hệ thống chưa được kiểm thử với dữ liệu thật và người dùng thật trong môi trường y tế. Các luồng chức năng đã được kiểm tra bằng tài khoản mẫu và dữ liệu thử nghiệm, nhưng điều đó chưa đủ để đánh giá toàn diện hiệu quả sử dụng trong thực tế. Nếu triển khai tại một phòng khám hoặc bệnh viện, hệ thống cần được kiểm thử với bác sĩ, y tá và bệnh nhân thật để đánh giá các yếu tố như mức độ dễ sử dụng, tần suất nhập liệu, cách phản hồi cảnh báo, số lượng cảnh báo phát sinh, thời gian bác sĩ cần để theo dõi bệnh nhân và mức độ phù hợp với quy trình làm việc hiện có.

\par Vấn đề thứ tư liên quan đến bảo mật và phân quyền dữ liệu. Hệ thống đã có các vai trò bệnh nhân, bác sĩ, y tá và quản trị viên, nhưng trong tương lai vẫn cần tiếp tục rà soát kỹ hơn các API để đảm bảo mỗi người dùng chỉ truy cập đúng dữ liệu được phép. Đặc biệt, bác sĩ và y tá không chỉ cần được kiểm tra theo vai trò, mà còn cần được kiểm tra theo quan hệ phân công với bệnh nhân. Ngoài ra, hệ thống cần tiếp tục hoàn thiện các vấn đề như thời gian sống của token, bảo vệ refresh token, mã hóa dữ liệu nhạy cảm, quản lý secrets, ghi nhận audit log đầy đủ và kiểm soát truy cập ở môi trường production.

\par Vấn đề thứ năm là hạ tầng triển khai hiện tại phù hợp với mục tiêu demo và kiểm thử đồ án, nhưng chưa phải mô hình tối ưu cho triển khai quy mô lớn. Backend hiện được triển khai trên Amazon EC2 bằng Docker Compose, các service như API server, worker, Redis, Temporal và Nginx chạy trong cùng một môi trường. Mô hình này đơn giản, dễ quản lý và phù hợp với giai đoạn thử nghiệm, nhưng nếu số lượng người dùng tăng, hệ thống cần được mở rộng về khả năng chịu tải, sao lưu dữ liệu, giám sát lỗi, tự động khôi phục, cân bằng tải và tách biệt các thành phần quan trọng.

\par Vấn đề thứ sáu là trải nghiệm người dùng trên mobile và web vẫn có thể được cải thiện thêm. Đối với mobile, các màn hình nhập chỉ số cần được tối ưu hơn cho bệnh nhân lớn tuổi, các thiết bị có kích thước màn hình khác nhau và các trường hợp mạng không ổn định. Đối với web bác sĩ, cần tiếp tục cải thiện cách ưu tiên cảnh báo, cách lọc dữ liệu, cách hiển thị biểu đồ và cách giúp bác sĩ nhanh chóng nhận biết bệnh nhân nào cần được theo dõi trước. Đối với trang quản trị, cần bổ sung thêm các cơ chế kiểm tra dữ liệu đầu vào, xác nhận thao tác nguy hiểm và phân loại lịch sử hoạt động rõ ràng hơn.

\par Cuối cùng, hệ thống hiện chưa tích hợp với các hệ thống hồ sơ bệnh án điện tử hoặc quy trình quản lý bệnh nhân có sẵn tại cơ sở y tế. Trong thực tế, nếu một hệ thống RPM được sử dụng tại bệnh viện hoặc phòng khám, dữ liệu bệnh nhân có thể cần đồng bộ với các hệ thống khác như hồ sơ bệnh án, lịch hẹn, đơn thuốc hoặc hệ thống quản lý nhân sự y tế. Việc chưa có tích hợp này giúp phạm vi đồ án gọn hơn, nhưng cũng là một điểm cần xem xét nếu muốn đưa hệ thống vào sử dụng thực tế.

\section{Hướng phát triển trong tương lai}
\par Ngoài việc tiếp tục khắc phục những vấn đề còn tồn đọng, nhóm cũng đề xuất một số hướng phát triển nhằm hoàn thiện hệ thống theo dõi bệnh nhân từ xa trong tương lai. Những hướng phát triển này tập trung vào việc nâng cao độ tin cậy của dữ liệu, mở rộng chức năng y tế, cải thiện trải nghiệm người dùng, tăng cường bảo mật và nâng cấp hạ tầng vận hành.

\begin{itemize}
    \item Tích hợp trực tiếp với thiết bị đo sức khỏe: Trong tương lai, hệ thống có thể được mở rộng để kết nối với máy đo huyết áp, máy đo đường huyết, cân thông minh, đồng hồ thông minh hoặc các thiết bị theo dõi sức khỏe khác. Khi dữ liệu được đồng bộ tự động từ thiết bị, hệ thống có thể giảm sai sót nhập liệu thủ công, tăng tần suất ghi nhận dữ liệu và giúp bác sĩ có cái nhìn chính xác hơn về diễn biến sức khỏe của bệnh nhân. Việc tích hợp có thể được thực hiện thông qua Bluetooth, API của nhà sản xuất thiết bị hoặc các nền tảng trung gian như Apple Health và Google Fit nếu phù hợp.

    \item Cá nhân hóa ngưỡng cảnh báo và kế hoạch theo dõi: Hệ thống hiện đã hỗ trợ bác sĩ cấu hình ngưỡng cho từng bệnh nhân, nhưng trong tương lai có thể mở rộng thêm khả năng quản lý kế hoạch theo dõi cá nhân hóa. Mỗi bệnh nhân có thể có bộ chỉ số cần theo dõi, tần suất đo, thời điểm đo và ngưỡng cảnh báo khác nhau tùy theo bệnh lý, tuổi, mức độ nguy cơ và chỉ định của bác sĩ. Việc cá nhân hóa này sẽ giúp hệ thống phù hợp hơn với thực tế lâm sàng thay vì áp dụng một cách theo dõi giống nhau cho mọi bệnh nhân.

    \item Bổ sung các thuật toán phân tích dữ liệu và hỗ trợ ra quyết định: Khi hệ thống có đủ dữ liệu lịch sử, nhóm có thể nghiên cứu các phương pháp phân tích xu hướng, phát hiện bất thường, phân tầng nguy cơ hoặc dự đoán khả năng xuất hiện biến cố sức khỏe. Tuy nhiên, các chức năng này cần được triển khai thận trọng, có sự tham gia của chuyên gia y tế và cần được kiểm chứng trên dữ liệu thực tế. Trong giai đoạn đầu, hệ thống có thể chỉ dừng ở mức gợi ý hoặc hỗ trợ bác sĩ ưu tiên bệnh nhân cần theo dõi, thay vì đưa ra quyết định điều trị tự động.

    \item Cải thiện quy trình xử lý cảnh báo: Cảnh báo có thể được mở rộng theo hướng có trạng thái xử lý rõ ràng hơn, ví dụ chưa xem, đã xem, đang xử lý, đã liên hệ bệnh nhân, đã yêu cầu đo lại hoặc đã chuyển tuyến. Hệ thống cũng có thể bổ sung cơ chế phân loại mức độ ưu tiên, nhóm cảnh báo trùng lặp, tạm ẩn cảnh báo ít nghiêm trọng hoặc tự động nhắc bác sĩ nếu cảnh báo quan trọng chưa được xử lý trong một khoảng thời gian nhất định. Điều này giúp giảm tình trạng quá tải cảnh báo và giúp nhân viên y tế tập trung vào các trường hợp quan trọng.

    \item Bổ sung chức năng quản lý thuốc và tuân thủ điều trị: Đối với bệnh nhân tăng huyết áp và đái tháo đường, việc dùng thuốc đúng giờ và đúng liều đóng vai trò rất quan trọng. Trong tương lai, hệ thống có thể bổ sung module quản lý thuốc, cho phép bác sĩ tạo lịch dùng thuốc, bệnh nhân xác nhận đã dùng thuốc và hệ thống thống kê mức độ tuân thủ. Kết hợp dữ liệu dùng thuốc với dữ liệu huyết áp hoặc đường huyết có thể giúp bác sĩ đánh giá tốt hơn hiệu quả của kế hoạch điều trị.

    \item Bổ sung hướng dẫn sử dụng và nội dung giáo dục sức khỏe: Hệ thống có thể thêm phần hướng dẫn cho người dùng mới, giải thích cách nhập chỉ số, cách đo huyết áp đúng, cách đo đường huyết đúng thời điểm và cách hiểu các thông báo cảnh báo trong ứng dụng. Ngoài ra, hệ thống có thể cung cấp các nội dung giáo dục sức khỏe ngắn gọn về chế độ ăn uống, vận động, theo dõi bệnh mãn tính và các dấu hiệu cần liên hệ với nhân viên y tế. Những nội dung này giúp bệnh nhân sử dụng hệ thống hiệu quả hơn và chủ động hơn trong quá trình chăm sóc sức khỏe.

    \item Hoàn thiện bảo mật và quyền riêng tư dữ liệu: Vì hệ thống xử lý dữ liệu sức khỏe và thông tin cá nhân, bảo mật cần được ưu tiên trong các giai đoạn phát triển tiếp theo. Nhóm có thể bổ sung các cơ chế như rút ngắn thời gian sống của access token, quản lý refresh token an toàn hơn, mã hóa dữ liệu nhạy cảm, phân quyền chi tiết theo quan hệ phân công, audit log đầy đủ, kiểm tra bảo mật định kỳ, quản lý secrets bằng công cụ chuyên dụng và xây dựng chính sách lưu trữ dữ liệu rõ ràng. Đây là điều kiện quan trọng nếu hệ thống muốn tiến gần hơn đến khả năng triển khai thực tế.

    \item Nâng cấp hạ tầng triển khai trên AWS: Hạ tầng hiện tại có thể được cải thiện bằng cách bổ sung CloudWatch cho logging và monitoring, cơ chế backup định kỳ cho cơ sở dữ liệu và cảnh báo khi service gặp lỗi. Đối với frontend, Cloudflare hiện đã hỗ trợ HTTPS phía trước S3; trong tương lai, nhóm có thể tối ưu thêm caching hoặc giám sát hiệu năng phân phối nội dung nếu quy mô truy cập tăng. Nếu quy mô backend tăng, hệ thống có thể được chuyển từ mô hình Docker Compose trên một EC2 sang các dịch vụ container chuyên dụng hơn như Amazon ECS hoặc Kubernetes, đồng thời tách riêng API server, worker, Redis và Temporal để tăng khả năng mở rộng.

    \item Bổ sung kiểm thử tự động và kiểm thử hiệu năng: Trong tương lai, nhóm có thể bổ sung unit test cho backend, integration test cho API, end-to-end test cho các luồng chính và kiểm thử giao diện cho frontend/mobile. Ngoài ra, hệ thống cần được kiểm thử hiệu năng để đánh giá khả năng xử lý khi có nhiều bệnh nhân nhập dữ liệu, nhiều cảnh báo phát sinh hoặc nhiều bác sĩ truy cập dashboard cùng lúc. Việc có bộ kiểm thử tự động giúp giảm rủi ro khi phát triển thêm chức năng mới và giúp hệ thống ổn định hơn sau mỗi lần cập nhật.

    \item Triển khai thử nghiệm với người dùng thật: Một hướng phát triển quan trọng là triển khai thử nghiệm hệ thống với một nhóm nhỏ bệnh nhân, bác sĩ và y tá để thu thập phản hồi thực tế. Thông qua quá trình thử nghiệm, nhóm có thể đánh giá xem bệnh nhân có dễ nhập chỉ số hay không, bác sĩ có thấy cảnh báo hữu ích hay không, giao diện có gây nhầm lẫn hay không và quy trình theo dõi có phù hợp với thực tế làm việc hay không. Các phản hồi này sẽ là cơ sở quan trọng để cải thiện sản phẩm trước khi mở rộng phạm vi sử dụng.

    \item Tích hợp với các hệ thống y tế khác: Nếu hệ thống được phát triển xa hơn, nhóm có thể nghiên cứu khả năng tích hợp với hồ sơ bệnh án điện tử, hệ thống quản lý lịch hẹn, hệ thống đơn thuốc hoặc các hệ thống quản lý bệnh nhân tại phòng khám và bệnh viện. Việc tích hợp này giúp dữ liệu theo dõi từ xa không bị tách biệt khỏi quy trình chăm sóc chính, đồng thời giúp bác sĩ có cái nhìn đầy đủ hơn khi đánh giá tình trạng bệnh nhân.
\end{itemize}

\par Tóm lại, đề tài đã giúp nhóm xây dựng được một hệ thống theo dõi bệnh nhân từ xa có đầy đủ các thành phần nền tảng, từ ứng dụng di động, ứng dụng web, trang quản trị, backend đến hạ tầng triển khai. Mặc dù hệ thống vẫn còn một số hạn chế và cần tiếp tục hoàn thiện nếu triển khai thực tế, kết quả đạt được đã chứng minh tính khả thi của mô hình Remote Patient Monitoring trong việc hỗ trợ quản lý bệnh mãn tính. Đây là nền tảng quan trọng để nhóm tiếp tục mở rộng hệ thống theo hướng tự động hóa dữ liệu, cá nhân hóa theo dõi, nâng cao bảo mật và ứng dụng sâu hơn vào lĩnh vực y tế số.